% Teil II: Analogie-Tabelle und Gegenüberstellung.

\subsection*{Gegenüberstellung}
Zentrale Entsprechungen (Skalarform, Rotation um feste Achse):

\begin{center}
\renewcommand{\arraystretch}{1.35}
\begin{tabular}{lcc}
\toprule
\textbf{Grösse} & \textbf{Translation} & \textbf{Rotation}\\
\midrule
Kinematik-Grösse & $x,\ v,\ a$ & $\varphi,\ \omega,\ \alpha$\\
Trägheit & $m$ & $I$\\
Ursache & Kraft $F$ & Drehmoment $M$\\
Bewegungsgleichung & $F=ma$ & $M=I\alpha$\\
Impuls/Drehimpuls & $p=mv$ & $L=I\omega$\\
Grundgesetz & $F=\dot p$ & $M=\dot L$\\
kinetische Energie & $E=\frac12 m v^2$ & $E_\mathrm{rot}=\frac12 I\omega^2$\\
Arbeit & $W=\int F\,dx$ & $W=\int M\,d\varphi$\\
\bottomrule
\end{tabular}
\end{center}

\textbf{Bedeutung der Analogie:} Das Eiskunstläufer-Problem (Drehimpuls
erhalten, $I$ ändert sich) entspricht dem linearen Problem "Anzug, bei dem
sich die Masse ändert, aber $p=mv$ konstant bleibt"; die Struktur der
Lösungen ist identisch.

\begin{ergebnis}
Jeder Translationsgrösse steht eine Rotationsgrösse gegenüber; die
Gesetze ($F=ma$ vs.\ $M=I\alpha$, Impuls vs.\ Drehimpuls, Energie- und
Arbeitsdefinitionen) haben dieselbe mathematische Form.
\end{ergebnis}

\begin{kontrolle}
\emph{Konsistenz:} $M=I\alpha$ aus $\dot L=\frac{d}{dt}(I\omega)=I\dot\omega$
für konstantes $I$ - genau wie $F=m\dot v$. \checkmark
\emph{Dimensionen:} $M=[kg\,m^2/s^2]$, $I\alpha=[kg\,m^2/s^2]$. \checkmark
\end{kontrolle}

\begin{merksatz}
Wenn man eine rotierende Grösse erkennt, hat man die Lösung schon: Man
ersetzt $m\to I$, $v\to\omega$, $F\to M$, $p\to L$ und erhält
Bewegungsgleichung, Erhaltungssatz und Energie aus dem schon verstandenen
Translationsfall.
\end{merksatz}